\section{Realizing the scheduler by blocking inequalities}\label{sec:layers}

At a probe for bucket \(P\), the scheduler of
Theorem~\ref{thm:scheduler} must make enough of the corresponding red
rows tight to serve the client program. Each red row will be a lower
bound on an affine expression, called an \emph{output}. A module has two
ways to supply a tight output. It can attain one of its own bounds, or
it can transmit tight bounds from all of its children at one label.
Section~\ref{sec:vertex} proves this alternative for the inequalities
constructed here. Section~\ref{sec:compiler} chooses modules to match
the alternative to the options in the client program.

\paragraph{One episode.}
Write \(x_P\) for the probe of bucket \(P\) in a single episode.
One coordinate \(v_P\) per bucket suffices to realize the first alternative
at height zero. Impose
\begin{equation}\label{eq:base-module}
 v_P\ge0,\qquad v_P+v_{P'}\ge1\quad(P\ne P'),\qquad
 v_P\ge I_0|x-x_P|-1 .
\end{equation}
Choose \(I_0\) so that the time-dependent bound for \(P\) is at
least two at every other probe and at both anchors. At \(x_P\) its
value is \(-1\), while each retained \(v_{P'}\) for \(P'\ne P\) is at
least two. Thus \(v_P\) can move in both directions if it is present
and positive. At a vertex that retains it, the floor \(v_P=0\) must
be tight. Reusing a coordinate at two separated probes poses a new
problem. A convex lower bound low at those probes stays low between
them. The numerical construction provides many different zero sets
on the same coordinates. They will distinguish the sites in a time
block, while the prices introduced below select children to serve at
a probe.

\subsection{From ordered layers to prescribed times}

Fix \(r\ge100\). The positions \(\mathcal A\) will carry a tournament
in which each set of at most \(2^r\) positions has a common dominator
outside the set. Such a tournament exists. For a fixed \(j\)-set in
a uniform random tournament on \(a_0\) positions, with \(1\le j\le
2^r\), the probability of having no common dominator is
\((1-2^{-j})^{a_0-j}\). Multiplying by \(\binom{a_0}{j}\) and summing
over these \(j\) gives a bound tending to zero as \(a_0\to\infty\).
Fix the tournament once for this \(r\).

Apply Theorem~\ref{thm:coarse} with \(C=1000r^2\). Choose a power of
two \(q\) supplied by that theorem and take \(L=(q^C)^k\) for sufficiently
large \(k\). Its field size \(b\), the least power of two at least
\(L^{1/(r+1)}\), then accommodates all the positions. Every block
will use the same label alphabet
\[
   \mathcal I\times\mathcal A\times\mathbb F_b,\qquad
   H_i(s)=\{(i,a,\mathrm h_a(s)):a\in\mathcal A\}
    \quad(i\in\mathcal I,\ 1\le s\le L).
\]
Here \(H_i(s)\) is the bundle assigned to site \(s\) by index \(i\).
Since \(\lvert\mathcal I\rvert\le L^{8/C}\), the total number \(V\)
of labels satisfies
\[
 V\le 2|\mathcal A| L^{8/C+1/(r+1)},\qquad 10V\le L^{3/r}
\]
for all sufficiently large \(k\).

At each site the theorem allows a blank sample, any positive number of
interior samples, and another blank sample. At an interior sample of
site \(s\), every state vanishes on some whole \(H_i(s)\) and has at
most \(r\) zeros on any bundle of a different site. Its third conclusion
bounds the number of states at a sample independently of the number of
interior samples. We now use its merge inequality to realize these
states at the \emph{prescribed real times}, including unevenly spaced ones.

For a sample time \(x\) with immediate neighbors \(x^-<x<x^+\), write
\(x=\alpha_x x^-+(1-\alpha_x)x^+\). Choose a single positive integer
\(e\) for the finite sequence so that \(4^{-e}<\min(\alpha_x,1-\alpha_x)\)
at every such \(x\). Raise each entry of every state to the power \(e\).
The zero sets persist and the merge factor in
Equation~\eqref{eq:layermerge} becomes \(4^{-e}\). Let \(Z_x\) be the
convex hull of the powered states at \(x\) plus the nonnegative
orthant. For \(c\ge0\), the value
\(m_c(x)=\min\{c\cdot v:v\in Z_x\}\) is attained by a powered state.
Merging minimizing states at the two neighbors gives
\[
 m_c(x)\le 4^{-e}\bigl(m_c(x^-)+m_c(x^+)\bigr)
        \le \alpha_x m_c(x^-)+(1-\alpha_x)m_c(x^+).
\]
The last inequality is strict whenever \(m_c(x)>0\), since the two
neighboring values then have positive sum. Thus the linear interpolation
of the data \(m_c\) is convex and bends strictly at each interior time
where the data are positive. At any sample \(x_0\) with \(m_c(x_0)>0\),
choose an affine function \(b_c\) equal to \(m_c\) there and strictly
below it at every other sample. At an interior sample one may use any
slope strictly between the two chord slopes. At the first or last
sample use a slope below or above its adjacent chord slope, respectively.

\subsection{Time blocks and the forest of modules}
Time blocks organize the sample times. Modules organize the physical
coordinates. Blocks at one height are disjoint intervals of time,
and the same module serves all of them.

Arrange the $N=L^h$ episodes as the leaves of an ordered $L$-ary
tree of time blocks. A block of height zero contains the anchors
and probes of one episode. A block $E$ of positive height has one
site $s$ for each of its children $E_s$. All samples of $E_s$,
including its own blanks if any, are interior samples of that site.
Insert two new blank samples around them. The gaps between prescribed
episodes leave room for all internal insertions, and the outermost
blanks can be placed beyond the episodes. Every prescribed probe is an
interior sample of its site at each positive height, so the zero-bundle
conclusion applies there. The added blanks serve only the interpolation
within the block that inserted them. A common exponent \(e\) can be chosen
for the finitely many interpolation triples in all the blocks. The
total number of distinct times, including the new blanks, is
$S=O(NB)$, since two blanks are added per tree edge. Whenever sample
data on the whole tree are interpolated, we use linear interpolation in
their order and extend the two outermost pieces affinely to the line.

A separate finite forest specifies which coordinates we use. Each
module owns distinct coordinates. The roots are
modules of height $h$. A module of height $t>0$ may have any number of
children of height $t-1$ at each label, and modules of height zero
have no children. Its own outputs, called quick outputs, come in classes with one member
per bucket $P\in\{1,\ldots,B\}$. It also propagates each output class
of each child as a separate class of its own. Thus a root has one
class for each quick-output class in its subtree. The rows bounding
these root outputs below will be the red rows. At any fixed time the
root systems use disjoint coordinates, so their joint slice is a product.

All modules use one physical time $x$. Each module has a fixed set of
physical coordinates shared across its blocks. To state the inequalities
for block $E$, take the local value of every output to be its physical
value minus
\begin{equation}\label{eq:module-shift}
 D_E(x)=T_E(x)+R_t(x),
\end{equation}
where $T_E$ is affine and $R_t$ is convex piecewise affine, common
to all blocks of this height. At height $h$ we set the shift to zero.
At the samples in \(E\), its lines and shifts will make the local system
of this block govern its module. We first give those local inequalities.

At height zero there is one physical coordinate per bucket. Subtracting
the shift \eqref{eq:module-shift} gives the local coordinates $v_P$ in
\eqref{eq:base-module}. Choose a common $I_0$ large enough to serve all
height-zero blocks.

\subsection{Local inequalities at positive height}
Fix $t>0$ and a block $E$ of height $t$.
The child systems are already fixed. Their pair bounds have positive
gaps, and Lemma~\ref{lem:descent} at height \(t-1\) supplies an allowance
for raising their local output floors independently at probes. Choose
\(\delta>0\) with \(10\delta\) smaller than all those gaps and with
\(2\delta\) within the child floor allowance. In assigning the shifts
at height \(t-1\), take \(0\le R_{t-1}\le\delta/8\) at all samples.
The final paragraph describes the choices in inductive order.

At each label $\ell$ introduce local variables $W_\ell,p_\ell$ and,
for every bucket $P$, quick outputs $F_{\ell,P},G_{\ell,P}$. The
$F$'s at a fixed label constitute one class, and the $G$'s another.
If a child attached at $\ell$ has physical output $v_j^{\rm phys}$,
its propagated output in this block is the expression
\begin{equation}\label{eq:module-propagation}
 Y_j=p_\ell+\varepsilon v_j^{\rm phys}.
\end{equation}
The index $j$ includes the child identity. No new coordinate is
introduced for $Y_j$. Take $\mu=\sigma^2$ and $\varepsilon=\sigma^3$,
where $\sigma>0$ will be chosen small at this height. Impose output
floors $F,G,Y\ge0$. A floor raise in a fixed slice refers to an
increase of one of these local lower bounds.

\paragraph{Numerical constraints.}
Use the states and $Z_x$ for this block, and require $W\ge0$.
Choose a piecewise affine function $g$ with positive values at
all samples, so that adding any offsets in $\{0,1\}$ to its sample
values still gives concave data. A sufficiently large multiple of
$-x^2$ on the finite sample set, translated upward and interpolated,
has this property. We use the same \(g\) on the full sample set at this
height. Write $z=W-\sigma g(x)\mathbf1$ and let $z^+,z^-$
denote its positive and negative parts, both taken nonnegative.
The numerical rows must reproduce the following condition at each
sample \(x\) of \(E\):
\begin{equation}\label{eq:layer-state}
 z^+\in D Z_x,\qquad D=1+M_0\sum_\ell z_\ell^- .
\end{equation}
A coordinate of \(W\) is allowed to fall below the offset
\(\sigma g(x)\). When that happens the negative parts of \(z\) enlarge
the required positive part by the factor \(D\). The following rows
implement this compensation with nonnegative coefficients.

Apart from nonnegativity, $Z_{x_0}$ is described by its facet
inequalities $c\cdot v\ge m$, with $c\ge0$ and $m=m_c(x_0)>0$.
Indeed its recession cone is the nonnegative orthant, and an
inequality with minimum zero already follows from $v\ge0$.
The number of facets is bounded in the fixed parameters, since each
face is generated by a subset of the boundedly many states and coordinate
rays. For all the finitely many facets and samples in use take $M_0$
large enough that $M_0m\ge\max_\ell c_\ell$. If $b_c$ is the affine
minorant of $m_c$ touching only at $x_0$, impose
\begin{equation}\label{eq:coarse-rows}
 c'\cdot W\ge b_c(x)+\sigma\|c'\|_1g(x_0)
\end{equation}
for each $c'$ formed by replacing any subset of the entries of $c$
by $M_0m$. At $x_0$ the least value of $c'\cdot z$ over these
choices is $c\cdot z^+-M_0m\sum_\ell z_\ell^-$. Hence the rows
generated there enforce exactly \eqref{eq:layer-state}. At any other
sample $x$ of $E$, its own rows suffice. Since $W\ge z^+$ and $c'\ge c$,
we have $c'\cdot W\ge m_c(x)$ by \eqref{eq:layer-state} and the
inclusion $D Z_x\subseteq Z_x$ for $D\ge1$. The strict gap
$m_c(x)-b_c(x)$ makes \eqref{eq:coarse-rows} strict there for small
enough $\sigma$.

\paragraph{Links and prices.}
Let $o_P$ have value one at probes for $P$ and zero at every other
sample. At such a probe the quick outputs in $P$ will be called
\emph{target outputs}, and those in the remaining buckets \emph{other-bucket outputs}. Impose
\begin{equation}\label{eq:module-link}
 F_{\ell,P}\ge W_\ell-\sigma (g+o_P)(x),\qquad
 p_\ell\ge W_\ell-\sigma g(x)-3\mu .
\end{equation}
At the target probe the first link allows one extra \(\sigma\) of slack.
In Lemma~\ref{lem:descent}, this slack frees the price from its link to
\(W\) at a zero coordinate of the limiting state. A tournament on each
bundle will select one of the freed prices for the common descent. Enforce $p_\ell\ge-K_0\varepsilon$, with $K_0$ so large
that $K_0-T_{E_s}(x)-\delta>0$ for all relevant sites and samples.
For every site $s$ set
\[
 u_\ell^s=-p_\ell/\varepsilon-T_{E_s}(x)-\delta .
\]
Thus \(u_\ell^s\) measures the price in the units of a child at site
\(s\), with sign reversed. The price floor gives a common upper bound
\(K_1>0\) for these quantities at the samples. Choose $M_1>1$ such that $rK_1/M_1<\delta/2$.
For every $H_i(s)$ and every list $\ell_0,\ldots,\ell_r$ in it whose
positions form a transitive tournament in the stated order (each
beats all later positions), impose
\begin{equation}\label{eq:module-tournament}
 \sum_{j=0}^r M_1^{r-j}u_{\ell_j}^s\le0 .
\end{equation}

For the link to the second quick output put
$\mathcal D(x)=\max_{E'}D_{E'}(x)$ over all blocks of height $t-1$.
Take $0<\lambda<\delta/4$ so small that the sample values of
$\mathcal D+\lambda-2\lambda o_P$ are convex for every $P$.
This is possible because $\mathcal D=R_{t-1}+\max T_{E'}$ has strict
bends in its sample data. With the expression in parentheses
interpolated from those values, impose
\begin{equation}\label{eq:module-g}
 G_{\ell,P}\ge p_\ell+\varepsilon
       (\mathcal D+\lambda-2\lambda o_P)(x).
\end{equation}
Finally give each pair of distinct outputs within a class a positive
gap. For their sum the lower bound is
\begin{equation}\label{eq:module-pairs}
 \mu\quad\text{in an $F$ class},\qquad
 \varepsilon\lambda/4\quad\text{in a $G$ class},\qquad
 \varepsilon\delta\quad\text{in a propagated class}.
\end{equation}

\subsection{One blocking system for all blocks}
At positive height the actual coordinates of a module are
\begin{equation}\label{eq:module-coordinates}
 W_\ell^{\rm phys},\quad p_\ell^{\rm phys}-W_\ell^{\rm phys},\quad
 F_{\ell,P}^{\rm phys}-W_\ell^{\rm phys},\quad
 G_{\ell,P}^{\rm phys}-p_\ell^{\rm phys}.
\end{equation}
We refer to these, in order, as the coordinates associated to $W,p,F,G$.
For block $E$, subtract $D_E$ from the physical variables $p,F,G$ and
$D_E/2$ from $W$ to obtain local variables. Descendant variables in
\eqref{eq:module-propagation} keep their physical values. This makes
the physical propagated output $p_\ell^{\rm phys}+\varepsilon
v_j^{\rm phys}$, with the same shift as a quick output.

This choice of coordinates separates the differences in the links from
the sums in output floors. For example, \(p_\ell^{\rm phys}\) is the sum
of the first two coordinates, and \(G_{\ell,P}^{\rm phys}\) is their sum
plus the fourth. Propagation adds a positive multiple of a descendant
output. Thus each output is a nonnegative linear combination of blocking
coordinates. We verify the convexity of the corresponding time bounds.
In particular, \eqref{eq:module-link} and \eqref{eq:module-g} bound single
difference coordinates. The former acquires the convex shift $D_E/2$
and the latter has zero shift. The negative interpolants of \(g\) and
\(g+o_P\) are convex, as is the interpolant in the \(G\) link. The numerical rows acquire nonnegative
multiples of $D_E$. Equation~\eqref{eq:module-tournament} bounds a
positive weighted sum of physical prices below by an affine function
plus a positive multiple of $D_E$. Output floors and pair rows also
acquire nonnegative multiples of $D_E$. Each such lower bound
represents the finite list of its affine pieces, all imposed on the
whole time line. Even after removing every output floor, a bound on
each single coordinate remains, from $W\ge0$ and the links, or from
the absolute-value bound at height zero.

Choose the lines $T_E$ at each height with slopes increasing in the
order of the blocks and with successive intersections lying in the
gaps between them. Then $T_E$ is the unique highest line at every
sample of $E$. Scaling all these lines by a common positive factor
makes the gaps between the highest line and the others arbitrarily
large at those finitely many samples. In particular
$\mathcal D(x)=D_{E_s}(x)$ at the samples of $E_s$.
We check that the current block imposes all the required inequalities
at its samples. Express a feasible point for $E$ in the local variables of
another block $E'$. The values of $F,G,p$ and all propagated outputs
increase by $D_E-D_{E'}$, and $W$ increases by half that amount.
The difference $G_{\ell,P}-p_\ell$ is invariant under this change,
and its lower bound in \eqref{eq:module-g} is shared by all blocks at
height $t$. Thus the copies of this link in $E$ and $E'$ have exactly
the same slack, which may be zero. For each remaining row of $E'$
the left side gains a positive multiple of $D_E-D_{E'}$. Before this
gain its value is bounded below uniformly on the current slice.
Use the floors and links for individual quantities and differences,
$W\ge0$ for numerical rows, and the price floors for the tournament.
There are finitely many rows and samples. One scaling therefore gives
strict slack in every off-block row other than the shared links.
At height zero every off-block row gains positive slack by the same
argument. Raising current output floors preserves the implications. The conclusion also holds after coordinate restriction. The system of
\(E\) and its descendants has nonnegative coordinate normals, so any
feasible restricted point can be lifted by taking omitted coordinates
sufficiently large. The full lift satisfies the other blocks' rows,
and hence so does its projection for every such row that is retained.

The construction in height is inductive. First fix all times,
numerical states, and powers. Construct the local systems at height
zero. Their pair gaps and allowable raises of the output floors
determine how small the next $\delta$ must be. After each height
$t<h$ has acquired its local pair gaps and positive floor allowance,
choose the next $\delta$ accordingly, then choose $R_t$ by interpolating
a sufficiently small positive multiple of $x^2$ on the full sample
set, so that $0\le R_t\le\delta/8$ there. Choose the lines $T_E$ at
that height as above. Thus when the next height is constructed all
child shifts are fixed. At a positive height first choose
$g,M_0,K_0,K_1,M_1,\lambda$, then choose a common sufficiently small
$\sigma$ for all its modules and blocks. Lemma~\ref{lem:descent}
gives a uniform floor allowance in these local systems, using only
their probe slices and the already fixed child systems. Thus its
choice of $\sigma$ and allowance is independent of the current
height's shifts. At the root height there is one
block, its shift is zero, and its output floors are single affine rows.
